%{{{Preamble
\documentclass[a4paper,10pt]{article} 
\usepackage{anysize}
\marginsize{2cm}{2cm}{1cm}{1cm}
%\textwidth 6.0in \textheight = 664pt
\usepackage{xltxtra}
\usepackage{graphicx}
\usepackage{subfig}
\usepackage{color}
\usepackage{xgreek}
\usepackage{framed}
\usepackage{float}
\usepackage{tkz-graph}
\setmainfont[Mapping=TeX-text]{CMU Serif}

%\setmainfont{Arial}
%\setsansfont{FreeSans}
%\setmonofont{FreeMono}
\begin{document}
\def\thesubsection {\alph{subsection}}
%\renewcommand{\labelenumi}{\roman{enumi})}
\renewcommand{\labelenumi}{\alph{enumi}}
\renewcommand{\labelenumii}{(\arabic{enumii})}
%}}}
%{{{Title
\begin{titlepage}
\begin{center}
\begin{figure}[t] 
     \includegraphics[scale=0.7]{title/ntua_logo}
\end{figure}
\begin{LARGE}\textbf{ΕΘΝΙΚΟ ΜΕΤΣΟΒΙΟ ΠΟΛΥΤΕΧΝΕΙΟ\\}\end{LARGE}
\vspace{5cm}
\begin{Large}
ΣΧΟΛΗ ΗΜ\&ΜΥ\\
Αλγόριθμοι και Πολυπλοκότητα\\
1\textsuperscript{η} Σειρά Γραπτών Ασκήσεων\\
Ακ. έτος 2010-2011\\
\end{Large}
\vfill
\begin{flushright}
\Large{Λύρας Γρηγόρης}\\
\large{Α.Μ.: 03109687}\\
\end{flushright}
\large\today\\
\end{center}
\end{titlepage}


%}}}
%{{{ 1 Προβολή ταινιών
\section{Προβολή ταινιών} \setcounter{section}{1}

Κρατάμε για κάθε ταινία $M_i$ έναν ακέραιο $Delta_i$ τον οποίο αρχικοποιούμε
στο 0 καθώς και μια σημαία που δηλώνει αν αυτή η ταινία έχει επιλεγεί από
κάποιο συνδρομητή. Καθώς περνάμε τις προτιμήσεις του χρήστη $a$ σημειώνουμε
για την ταινία της επιλογής του τη σημαία της σε "ΕΠΙΛΕΓΜΕΝΗ" και αυξάνουμε
τον μετρητή της αν αυτή είναι σημειωμένη για προβολή Σαββάτου από τον χρήστη
ενώ μειώνουμε αν αυτή έχει επιλεγεί για Κυριακή. Αφού ολοκληρώσουμε για όλους
τους χρήστες. διαβάζουμε τον συντελεστή κάθε των ταινίας και αν αυτός είναι
θετικός την τοποθετούμε στο σύνολο του Σαββάτου αν είναι αρνητικός στο σύνολο
της Κυριακής. Αν είναι 0 και έχει σημαία "ΕΠΙΛΕΓΜΕΝΗ" δεν έχει σημασία πού θα
επιλέξουμε να την βάλουμε.  Αν η σημαία δεν είναι "ΕΠΙΛΕΓΜΕΝΗ" τότε αγνοούμε
αυτή την ταινία.
%}}}

%{{{ 2 Μέτρηση συντομότερων μονοπατιών
\section{Μέτρηση Συντομότερων Μονοπατιών}
Ξεκινάμε από το άκρο s του γράφου. Θέτουμε τον κόμβο "υπό εξερεύνηση" και
ακολουθούμε όσες ακμές εξέρχονται από αυτόν. Θέτουμε τον κόμβο ως
"εξερευνημένο" και προχωρούμε στους νέους κόμβους. Αν με αυτή τη διαδικασία
φτάσουμε σε κόμβο "εξερευνημένο" είτε "υπό εξερεύνηση" τότε τον αγνούμε. Μόλις
φτάσουμε στο άκρο t ολοκληρώνουμε το τρέχον επίπεδο και σταματάμε τη
διαδικασία. Ο αριθμός των συντομότερων μονοπατιών είναι ο αριθμός των t που
εμφανίζεται στην τελευταία λίστα.

Η χρονική πολυπλοκότητα του αλγορίθμου είναι $\Theta(n)$
%}}}

%{{{ 3 Ελάχιστο Συνδετικό Δέντρο Υπό Περιορισμούς
\section{Ελάχιστο Συνδετικό Δέντρο Υπό Περιορισμούς}
\subsection{}
\begin{figure}[ht]
	\centering
	\subfloat[Initial]{
		\begin{tikzpicture}

			\Vertex[x=0,y=0,L=1]{1}
			\Vertex[x=3,y=0,L=2]{2}
			\Vertex[x=0,y=-3,L=3]{3}
			\Vertex[x=3,y=-3,L=4]{4}

			\Edge[label=2](1)(2)%
			\Edge[label=8](1)(3)
			\Edge[label=4](2)(3)
			\Edge[label=3](3)(4)
		\end{tikzpicture}
	}
	\hspace{1cm}
	\subfloat[MST]{
		\begin{tikzpicture}
			\Vertex[x=0,y=0,L=1]{1}
			\Vertex[x=3,y=0,L=2]{2}
			\Vertex[x=0,y=-3,L=3]{3}
			\Vertex[x=3,y=-3,L=4]{4}

			\Edge[label=2](1)(2)%
			\Edge[label=4](2)(3)
			\Edge[label=3](3)(4)
		\end{tikzpicture}
	}
	\hspace{1cm}
	\subfloat[L=\{2,4\}]{
		\begin{tikzpicture}
			\Vertex[x=0,y=0,L=1]{1}
			\Vertex[x=3,y=0,L=2]{2}
			\Vertex[x=0,y=-3,L=3]{3}
			\Vertex[x=3,y=-3,L=4]{4}

			\Edge[label=2](1)(2)%
			\Edge[label=8](1)(3)
			\Edge[label=3](3)(4)
		\end{tikzpicture}
	}
\end{figure}

\subsection{}
Αφαιρούμε τις κορυφές του L και βρίσκουμε το ελάχιστο συνδετικό δέντρο χωρίς
αυτές. Τέλος προσθέτουμε τις κορυφές του L ως φύλλα στο δέντρο που
φτιάξαμε χρησιμοποιώντας τις πιο "φτηνές" ακμές τους.

%}}}

%{{{ 4 Μοναδικότητα Ελάχιστου Συνδετικού Δέντρου
\section{Μοναδικότητα Ελάχιστου Συνδετικού Δέντρου}
\subsection{}
\begin{figure}[ht]
	\centering
	\subfloat[Initial]{
		\begin{tikzpicture}

			\Vertex[x=0,y=0,L=1]{1}
			\Vertex[x=3,y=0,L=2]{2}
			\Vertex[x=0,y=-3,L=3]{3}
			\Vertex[x=3,y=-3,L=4]{4}
			\Vertex[x=6,y=-3,L=5]{5}

			\Edge[label=2](1)(2)%
			\Edge[label=8](1)(3)
			\Edge[label=4](2)(3)
			\Edge[label=3](3)(4)
			\Edge[label=3](4)(5)
		\end{tikzpicture}
	}
	\hspace{1cm}
	\subfloat[MST]{
		\begin{tikzpicture}
			\Vertex[x=0,y=0,L=1]{1}
			\Vertex[x=3,y=0,L=2]{2}
			\Vertex[x=0,y=-3,L=3]{3}
			\Vertex[x=3,y=-3,L=4]{4}
			\Vertex[x=6,y=-3,L=5]{5}

			\Edge[label=2](1)(2)%
			\Edge[label=4](2)(3)
			\Edge[label=3](3)(4)
			\Edge[label=3](4)(5)
		\end{tikzpicture}
	}
\end{figure}
\subsection{}
\begin{figure}[ht]
	\centering
	\subfloat[Initial]{
		\begin{tikzpicture}

			\Vertex[x=0,y=0,L=$S$]{S}
			\Vertex[x=3,y=0,L=$V \backslash S$]{VS}

			\Edge[label=a](S)(VS)%
			\Edge[label=b > a,style={bend right}](S)(VS)%
			\Edge[label=c > a,style={bend left}](S)(VS)%
		\end{tikzpicture}
	}
	\hspace{1cm}
	\subfloat[MST]{
		\begin{tikzpicture}

			\Vertex[x=0,y=0,L=$S$]{S}
			\Vertex[x=3,y=0,L=$V \backslash S$]{VS}

			\Edge[label=a](S)(VS)%
		\end{tikzpicture}
	}
\end{figure}
Με δεδομένο πως αυτή η ακμή ανάμεσα σε μια τομή θα ενώνει δύο υποδέντρα που 
πάλι θα έχουν ελάχιστο βάρος, αν η ακμή με το ελάχιστο βάρος είναι μοναδική,
τότε αυτή θα είναι σίγουρα στο τελικό MST. Και εφόσον γνωρίζουμε πως σε κάθε
τομή η ακμή με το ελάχιστο βάρος είναι μοναδική, τότε για κάθε τομή μπορούμε
να επιλέξουμε ακριβώς μία ακμή για να συνθέσουμε το MST. Συνεπώς το MST είναι
μοναδικό.


\begin{figure}[ht]
	\centering
	\subfloat[Initial]{
		\begin{tikzpicture}

			\Vertex[x=0,y=0,L=1]{1}
			\Vertex[x=3,y=0,L=2]{2}
			\Vertex[x=0,y=-3,L=3]{3}
			\Vertex[x=3,y=-3,L=4]{4}

			\Edge[label=2](1)(2)%
			\Edge[label=8](1)(3)
			\Edge[label=2](2)(3)
			\Edge[label=2](3)(4)
			\Edge[label=2](2)(4)
		\end{tikzpicture}
	}
	\hspace{1cm}
	\subfloat[MST]{
		\begin{tikzpicture}
			\Vertex[x=0,y=0,L=1]{1}
			\Vertex[x=3,y=0,L=2]{2}
			\Vertex[x=0,y=-3,L=3]{3}
			\Vertex[x=3,y=-3,L=4]{4}

			\Edge[label=2](1)(2)%
			\Edge[label=2](2)(3)
			\Edge[label=2](2)(4)
		\end{tikzpicture}
	}
\end{figure}

Παρατηρούμε πως οι ακμές (2)-(3) και (2)-(4) έχουν και οι δύο το ίδιο βάρος
στο γράφο. Οπότε θεωρώντας $S=\{1,2\}$ έχουμε δύο ακμές με ελάχιστο βάρος να
διασχίζουν την $(S, V \backslash S)$. Τελικά όμως το MST είναι μοναδικό.

\subsection{}
Αν σε κάθε τομή κάθε ακμή ελάχιστου βάρους ανήκει στο MST τότε αυτό είναι μοναδικό.

Έστω πως δεν είναι μοναδικό, τότε θα υπάρχει κάποια ακμή σε κάποια τομή η
οποία δε θα έχει το ελάχιστο βάρος και θα ανήκει στο MST. Αν ισχύει αυτό τότε
αντικαταστήσαμε μια ακμή $e_1$ με μιά ακμή $e_2$ μεγαλύτερου βάρους. Συνεπώς
το δέντρο πλέον δεν είναι ελάχιστου βάρους.

\subsection{}
Στην ουσία χρειάζεται να ελέγξουμε κατά πόσο όλες οι ελάχιστες ακμές κάθε
τομής ανήκουν στο MST. Για αυτό το σκοπό εργαζόμαστε ως εξής:
\begin{itemize}
	\item Ταξινομούμε τις ακμές με βάση το βάρος τους.
	\item Έχουμε ένα σύνολο κορυφών S στο οποίο θα δημιουργήσουμε το MST
	\item Τοποθετούμε μία μία τις ακμές από τη μικρότερη προς τη μεγαλύτερη
		χωρίς να δυμιουργήσουμε κύκλο.
	\item Αν δημιουργηθεί κύκλος και στην προηγούμενη επανάληψη βάλαμε ακμή με
		ίδιο βάρος σταματάμε, το MST δεν είναι μοναδικό.
	\item Αν ολοκληρώσουμε τη διαδικασία και δημιουργήσουμε το MST, τότε αυτό είναι
		μοναδικό.
\end{itemize}
%}}}

%{{{ 5 Υπολογισμός ελάχιστου συνδετικού δέντρου με διαγραφή ακμών
\section{Υπολογισμός Ελάχιστου Συνδετικού Δέντρου με Διαγραφή Ακμών}
\subsection{}
Έστω ότι δεν υπάρχει MST χωρίς την ακμή $e$. Έστω πως μία ακμή που λείπει από τον κύκλο
C η $e_1$. Τότε μπορούμε να τοποθετήσουμε την $e_1$ πάλι σε αυτό
δημιουργώντας ένα κύκλό C2. Συνεπώς μπορούμε να διαγράψουμε μια ακμή από το
γράφο για να ξαναγίνει δέντρο, και μάλιστα μπορούμε να διαγράψουμε την $e$
εφόσον η ακμή $e$ και η ακμή $e_1$ ανήκουν στον κύκλο που δημιουργήθηκε.
Συνεπώς διαγράφοντας την $e$ δημιουργούμε ένα καινούριο MST με μικρότερο
βάρος από το αρχικό συνεπώς το αρχικό δέντρο δεν ήταν ελάχιστο.

\subsection{}
Ο αλγόριθμος αυτός αφαιρεί τις βαρύτερες ακμές όπου είναι δυνατόν από το
γράφο αφήνοντας μόνο τις απαραίτητες ώστε να είναι συνδετικός, συνεπώς είναι ορθός.
\subsection{}
Για να εφαρμόσουμε τον αλγόριθμο πρέπει σε κάθε εισαγωγή να βρούμε αν η ακμή
που εξετάζουμε ανήκει σε κύκλο.

Δημιουργούμε ένα πίνακα όπου τοποθετούμε τους κύκλους που εμφανίζονται στο
γράφο καθώς και τις ακμές που τους αφορούν τους οποίους βρίσκουμε αρχικά 
με αναζήτηση κατά πλάτος.
Για κάθε ακμή χρησιμοποιούμε τον προηγούμενο πίνακα για να διαπιστώσουμε αν
ανήκει σε κάποιο κύκλο. Αν δεν βρούμε κάποιον τότε η ακμή παραμένει στο γράφο.
Αν ανήκει σε κύκλο την αφαιρούμε, διαγράφουμε τον κύκλο από τον πίνακα και
συνεχίζουμε στην επόμενη.
%}}}
\end{document}
